\section{Finite computations and reproducibility}\label{app:checks}
The code in \texttt{anc/audit/check\_models.py} represents an operator as a finite sum of weighted basis permutations. For each defining relation it enumerates all basis columns of the tensor factors that occur in that relation. The absent factors are identities, so averaging squared column norms computes exactly the full-space normalized Hilbert--Schmidt expression, before floating-point rounding. Equal output rows are combined before squaring, avoiding cancellation of large squared terms. The periodic coordinates in~\eqref{eq:series} are summed with integer numerators and denominator $4^\ell-1$ for the permutation period $\ell$, then rounded once to floating point.

The test program has one counter and five controls. From the starting control, zero halts, while nonzero decrements and moves to an odd-parity control. There zero enters an absorbing nonhalting control, while nonzero decrements and returns to start. Idle and halt are absorbing. Hence input $e$ halts exactly when it is even, at step $e+1$. This is a test of the finite-model mechanism, not a computation of the universal halting set.

\begin{table}[ht]
\centering
\caption{Finite-model checks for the disclosed test program. Each ratio is a flat Hilbert--Schmidt defect divided by $4^{-N}$. The first two rows apply separately to all three compulsory clocks.}\label{tab:checks}
\small
\begin{tabular}{lrrr}
\toprule
Quantity & $N=4$ & $N=6$ & $N=8$\\
\midrule
Clock affine ratio & $3/\sqrt2$ & $3/\sqrt2$ & $3/\sqrt2$\\
Clock tail-flip ratio & $3$ & $3$ & $3$\\
Guarded increment ratio & $0.353553$ & $0.176777$ & $0.088388$\\
Carry-step ratio & $0.25$ & $0.125$ & $0.0625$\\
Conditional tail ratio & $0.00585946$ & $0.000160228$ & $0.00256366$\\
\midrule
$K_N$ & $\{0\}$ & $\{0,2\}$ & $\{0,2\}$\\
$k_N$ & $2$ & $4$ & $4$\\
Dimension & $16\,384$ & $4\,194\,304$ & $268\,435\,456$\\
Recorded checks & $390$ & $394$ & $398$\\
\bottomrule
\end{tabular}
\end{table}

The remaining defining defects are zero to within $10^{-14}$ in these tests. Complete per-relation values and their reference scales are archived, not just the maxima in Table~\ref{tab:checks}. The two guarded increments have identical ratios in this example. The theorem's operator-norm bounds are proved in Section~\ref{sec:finite}; a script's maximum column norm is not generally an operator norm. The displayed dimensions are represented without dense matrices. The original $N=8$ run had an observed peak working set about $1.16$ GB.

The gadget script \texttt{anc/gadgets/gadget\_check.py} tests self-adjointness, squares, all pairwise commutations within each context, context products, and the auxiliary conjugacies on random finite instances. It is a numerical check of the explicit extension, not a proof of the converse; that converse is proved in Section~\ref{sec:compiler}.

The ancillary README gives commands, dependencies, recorded outputs and limitations. All computations are run with one numerical-library thread. The finite-model script caps its exhaustive workspace implementation at $N=8$ to stay within the stated memory budget. The independent engine described in the methods section is in the repository directory \texttt{verify/}, with its recorded output, and the two scripts of the independent proof check are in \texttt{anc/crosscheck/}.


\paragraph{Reader and extraction sections.} The spectrum program in \texttt{anc/spectrum/} builds 24 finite reader models at $N=2,4,6$ and eight rational thresholds, including both endpoints. It tests reset boundaries, reads beyond $N$, halt inclusion, slot collisions, dressing and defining operator residuals on the nontrivial kernels. Its other tests cover radical-inverse windows, compression with domain repairs, finite word moments and noncommuting-state Jensen. The protocol programs in \texttt{anc/extraction/} evaluate the stage inequality, full exact and perturbed protocols with mixed inputs and noncommuting controls, small perturbations with positive lower bounds, bounded adversarial searches, and one-sided inclusion repair. Their recorded outputs are preserved in the respective \texttt{recorded/} directories. The ancillary README states each file's scope and distinguishes the original allowances from the constants of Lemma~\ref{lem:inactive-extraction}.


\paragraph{Complementary-construction checks.} The script \texttt{anc/gate/toy\_gate.py} retains its original recorded results in \texttt{gate/recorded/} and writes a new result beside the script. It executes the signed packed polynomial instructions for three-colouring of $K_3$, a pair of noncommuting binary bases, and a flagged bad three-colouring of $K_4$. It checks contraction-register decompositions, initialized finite histories through every distinct reset period, cyclic boundaries, simultaneous data factors, decoder and reset recursions, complementary window counts and address dressing. This script uses a single projection family for both players; its history, reset and window tests are independent of that choice and cover 16,893 history transitions, 35 cyclic reset boundaries and 2,396,744 mask/start cases across 254 masks. The separate two-family suite in \texttt{anc/gate/crosscheck/} tests the presentation of Section~\ref{sec:gate-input}, with separate families for the two players: $K_3$ colouring has 18 generators, and the magic-square witness has 42 generators and packed dimension 168. It also checks the reduced-state eigenspaces, an anti-synchronous PCC control, 53,248 history points, sparse address models and windows through depth 22. Recorded logs and rerun commands are specified in the ancillary README. These are tests of the new matrix algebra and finite histories on known-answer toy games, not the $\mathrm{MIP}^*$ games themselves. They do not certify the halting reduction, arbitrary-state lower bound or compiled game.

\paragraph{Stop-or-continue classification.} The four scripts in \texttt{anc/classification/} test the new stop-or-continue algebra on finite known-answer instances. They cover oscillating rational approximations, finite stopping parameters, tail rejection, per-gate LOAD boundaries, matrix-valued registers, noncommuting witness tuples, reset periods, flagged default data and finite threshold/window budgets. The operator-level check evaluates the new equations as sparse operators built from a toy history, with negative controls for an omitted LOAD factor and imposed time invariance. The matrix-interface check displays the shifted-corner and stop-sector identities on a hand-built instance, in which the stop-sector identity holds by construction, with a negative control for dropped coordinate commutation. The toy scripts add negative controls that change data between the checks of a false gate (\texttt{toy\_classification.py}) or one step before a LOAD (\texttt{toy\_attack.py}). The two register checks initialize $M=I$ and $\Sigma=0$ on the outgoing LOAD edge. The operator-level script also checks a synthetic LOAD with incoming $M=0$, $\Sigma=I$, testing both the instruction history and the polynomial selectors against that transition; holding the incoming registers gives nonzero negative controls. Their recorded outputs are archived and hashed. They do not instantiate the $\mathrm{MIP}^*$ reduction, synthesize a universal counter compiler, certify a noncomputable limit, or test the full compiled game. The lower bound and the nonuniform slow-diagonal upper bound are written proofs in Section~\ref{sec:classification}, not conclusions of these finite tests.


\paragraph{Power-scale sections.} The seven programs in \texttt{anc/power/} are run, with fixed seeds, one process at a time and one numerical-library thread. The two moment scripts enumerate finite-field code moments, check full-Fock CAR and small affine actions, and rebuild the collective clock in natural and sparse coordinate representations, including $M=48$. They verify finite examples of the $27$-row certificate's geometric ingredients; they do not allocate the full large affine representation. The stopped and rotation scripts test finite cyclic histories, zero and dyadic inputs, reset interruptions and the inherited unguarded increment defect. The independent stack tests use small noncommuting matrices and diagonal histories to stress ordered updates, column completeness, polynomial carry/clock charges and scalar product envelopes. The scaled-machine script checks all new counter, control, witness and calibrator equations on finite clock blocks, literal noncommutative polynomial identities, weighted bad-zero-test perturbations and a finite scheduler arithmetic simulation. Its scheduler test is not a compiled universal register program; compilation and elementary-step bounds are proved in \cref{lem:power-stream}. The independently written counter tests use exact rational scalar histories, small weighted noncommuting finite-window perturbations and arbitrary ordered measurement columns; they do not include a complete vanishing-defect clock/compiler strategy. Their scope and every output are retained under the corresponding \texttt{recorded/} directories, with SHA-256 hashes. These checks support, and do not replace, \cref{thm:power-moments,prop:power-localization,prop:power-counter-stability,prop:power-calibration}.

\paragraph{Precision-profile checks.} The profile checks in \texttt{anc/profiles/} and their individual scopes are described here. The runner \texttt{anc/profiles/run\_checks.py} executes three suites. The carry suite uses exact rational diagonal/permutation calculations for the carry and unguarded overflow, reset and initialized/no-reset cylinders, Cantor Lipschitz constants and the checker-cost comparison. The hold suite simulates twelve million pulses, checks denominator padding and holds, exact calibrator histories for thirteen inputs, logarithmic profile accounting and power moment bands. Its scheduler is an abstract instruction-cost simulation, not a complete emitted universal machine table. The independent suite imports neither supplied engine; it tests overflow-state deletion through twelve bits, noncommuting shared inputs, exact greedy tracking, affine endpoint parameters, separate hold accounting and the diagonal bit decoder. Recorded JSON/logs are archived with SHA-256 hashes. Finite outcomes support these interfaces; no numerical run certifies a noncomputable endpoint, a uniform operator theorem or exclusion by every game.

\paragraph{Ceiling, pairwise and group checks.} Four suites in \texttt{anc/frontiers/} test the additional finite interfaces. The supplied ceiling program checks finite clock defects, exact scalar calibrator domains and residuals, sharp-versus-smooth penalties and stopped multiplier saturation. The supplied hash program exhaustively checks affine collisions and the pairwise dimension bound. Independently specified tests cover weighted sequential POVM positivity, telescoping and effect-approximation ratios, noncommuting complete columns, actual-failure neighborhoods, mixed classical hash flags and exact rational stopped histories. The group suite checks sixteen finite counter programs on joint fuel/time grids, the single-visit trace kernel including zero and powers of two, the seventeen-state squaring loop through four marker visits, and finite-Heisenberg matrix trace normalization. Its elementary marker executes 595,050 instructions after initialization. These finite checks do not prove the arithmetic-hierarchy theorem, asymptotic irrationality, arbitrary-state trace rigidity, or an algebraic group-transfer theorem. The written arguments give those conclusions. JSON and console outputs are archived and hashed; the runner uses one numerical-library thread.
